\documentclass[10pt, letterpaper]{article}

% Packages:
\usepackage[
	ignoreheadfoot, % set margins without considering header and footer
	top=1 cm, % seperation between body and page edge from the top
	bottom=1 cm, % seperation between body and page edge from the bottom
	left=1 cm, % seperation between body and page edge from the left
	right=1 cm, % seperation between body and page edge from the right
	footskip=1 cm, % seperation between body and footer
	% showframe % for debugging
]{geometry} % for adjusting page geometry
\usepackage{titlesec} % for customizing section titles
\usepackage[dvipsnames]{xcolor} % for coloring text
\definecolor{primaryColor}{RGB}{0, 0, 0} % define primary color
\usepackage{enumitem} % for customizing lists
\usepackage{etoolbox} % for conditional statements
\usepackage{amsmath} % for math
\usepackage[
	pdftitle={Curriculo de Marllon Rodrigues},
	pdfauthor={Marllon Rodrigues},
	pdfcreator={Marllon Rodrigues},
	colorlinks=true,
	urlcolor=primaryColor
]{hyperref} % for links, metadata and bookmarks
\usepackage[pscoord]{eso-pic} % for floating text on the page
\usepackage{calc} % for calculating lengths
\usepackage{bookmark} % for bookmarks
\usepackage{lastpage} % for getting the total number of pages
\usepackage{changepage} % for one column entries (adjustwidth environment)
\usepackage{paracol} % for two and three column entries
\usepackage{needspace} % for avoiding page brake right after the section title
\usepackage{iftex} % check if engine is pdflatex, xetex or luatex

% Ensure that generate pdf is machine readable/ATS parsable:
\ifPDFTeX
\input{glyphtounicode}
\pdfgentounicode=1
\usepackage[T1]{fontenc}
\usepackage[utf8]{inputenc}
\usepackage{lmodern}
\fi

\usepackage{charter}

% Some settings:
\raggedright
\AtBeginEnvironment{adjustwidth}{\partopsep0pt} % remove space before adjustwidth environment
\pagestyle{empty} % no header or footer
\setcounter{secnumdepth}{0} % no section numbering
\setlength{\parindent}{0pt} % no indentation
\setlength{\topskip}{0pt} % no top skip
\setlength{\columnsep}{0.15cm} % set column seperation
\pagenumbering{gobble} % no page numbering

\titleformat{\section}{\needspace{4\baselineskip}\bfseries\large}{}{0pt}{}[
\vspace{1pt}
\titlerule]

\titlespacing{\section}{
% left space:
-1pt }{
% top space:
0.3 cm }{
% bottom space:
0.2 cm } % section title spacing

\renewcommand{\labelitemi}{$\vcenter{\hbox{\small$\bullet$}}$} % custom bullet points
\newenvironment{highlights}{ \begin{itemize}[ topsep=0.10 cm, parsep=0.10 cm, partopsep=0pt,
itemsep=0pt, leftmargin=0 cm + 10pt ] }{ \end{itemize} } % new environment for highlights

\newenvironment{highlightsforbulletentries}{ \begin{itemize}[ topsep=0.10 cm,
parsep=0.10 cm, partopsep=0pt, itemsep=0pt, leftmargin=10pt ] }{ \end{itemize} } % new environment for highlights for bullet entries

\newenvironment{onecolentry}{ \begin{adjustwidth}{ 0 cm + 0.00001 cm }{ 0 cm + 0.00001 cm }
}{ \end{adjustwidth} } % new environment for one column entries

\newenvironment{twocolentry}[2][]{ \onecolentry \def\secondColumn{#2} \setcolumnwidth{\fill, 4.5 cm}
\begin{paracol}{2} }{ \switchcolumn \raggedleft \secondColumn \end{paracol}
\endonecolentry } % new environment for two column entries

\newenvironment{threecolentry}[3][]{ \onecolentry
\def\thirdColumn{#3} \setcolumnwidth{, \fill, 4.5 cm}
\begin{paracol}{3} {\raggedright #2} \switchcolumn }{ \switchcolumn \raggedleft \thirdColumn
\end{paracol} \endonecolentry } % new environment for three column entries

\newenvironment{header}{
\setlength{\topsep}{0pt}
\par\kern\topsep
\centering
\linespread{1.5} }{ \par\kern\topsep } % new environment for the header

\newcommand{\placelastupdatedtext}{% \placetextbox{<horizontal pos>}{<vertical pos>}{<stuff>}
\AddToShipoutPictureFG*{% Add <stuff> to current page foreground
\put( \LenToUnit{\paperwidth-2 cm-0 cm+0.05cm}, \LenToUnit{\paperheight-1.0 cm} ){\vtop{{\null}\makebox[0pt][c]{ \small\color{gray}\textit{Última atualização em Maio de 2026}\hspace{\widthof{Última atualização em Maio de 2026}} }}}%
}%
}%

% save the original href command in a new command:
\let\hrefWithoutArrow\href

% new command for external links:

\begin{document}
	\newcommand{\AND}{\unskip \cleaders\copy\ANDbox\hskip\wd\ANDbox \ignorespaces }
	\newsavebox{\ANDbox}
	\sbox{\ANDbox}{$|$}

	\begin{header}
		\fontsize{25 pt}{25 pt}\selectfont Marllon Rodrigues

        \fontsize{18 pt}{18 pt}\selectfont Engenheiro de Software \& Automação

		\vspace{5 pt}

		\normalsize
		\mbox{\hrefWithoutArrow{mailto:marllonrsantos@gmail.com}{marllonrsantos@gmail.com}}%
		\kern 5.0 pt%
		\AND%
		\kern 5.0 pt%
		\mbox{\hrefWithoutArrow{https://linkedin.com/in/rodriguesmarllon}{linkedin.com/in/rodriguesmarllon}}%
		\kern 5.0 pt%
        \AND%
		\kern 5.0 pt%
		\mbox{\hrefWithoutArrow{https://github.com/rodriguesmarllon}{github.com/rodriguesmarllon}}%
	\end{header}

	\section{Competências Técnicas}

    \textbf{Desenvolvimento de Software:} TypeScript, Node.js (Fastify, NestJS), Vue.js (Quasar), React, Next.js, Java, C\# (.NET), Python, Prisma, PostgreSQL, Redis, Cloudflare R2, Docker, CI/CD (GitHub Actions), SwiftUI, Git.

    \textbf{Automação Industrial:} PLC (Siemens S7-1200/1500, Rockwell, Schneider M580), HMI/SCADA (AVEVA Plant SCADA/Citect, System Platform, InTouch, Historian, FactoryTalk, ISA-101), Relés SEL, IEC 61850, Modbus TCP, Redes Industriais, Integração IT/OT.

	\section{Experiência}

	\begin{twocolentry}
		{ Set 2025 - Atual } \textbf{Co-Founder \& Tech Lead}, SIMP (Startup de Tecnologia Municipal)
	\end{twocolentry}

	\vspace{0.10 cm}
	\begin{onecolentry}
		\begin{highlights}
			\item \textbf{SaaS Multi-tenant (SIMP):} Arquitetura e desenvolvimento de um ERP de gestão municipal B2B do zero — multi-tenancy completo, RBAC, Kanban, notificações SSE, uploads Cloudflare R2. Stack: Fastify v5, React 19, Prisma, PostgreSQL, Redis; deploy Vercel + Render.
            \item \textbf{Inteligência Financeira:} Módulo de diagnóstico financeiro preditivo com Z-score e algoritmos de previsão ETS para detecção automática de anomalias em orçamentos municipais.
            \item \textbf{DevOps com IA:} Pipelines CI/CD com agentes Claude AI autônomos para revisão de código, análise de segurança (TruffleHog, CodeQL) e diagnóstico de falhas em cada commit e Pull Request.
			\item \textbf{Mentoria:} Capacitação de desenvolvedor júnior que evoluiu para atuação full stack autônoma, resolvendo issues e entregando features de forma independente.
		\end{highlights}
	\end{onecolentry}

    \vspace{0.2 cm}

    \begin{twocolentry}
		{ Jun 2025 - Atual } \textbf{Desenvolvedor Full Stack}, Bunkerchain (Singapura/Remoto)
	\end{twocolentry}

	\vspace{0.10 cm}
	\begin{onecolentry}
		\begin{highlights}
			\item \textbf{Plataforma eBDN (Methanol):} Desenvolvimento full stack no sistema de Nota de Entrega de Bunker eletrônica da Bunkerchain para operações de carga de Metanol — ferramenta crítica de conformidade marítima usada em tempo real.
            \item \textbf{Entrega de Feature (Quasar/Vue.js):} Desenvolvimento completo da funcionalidade Terminal Role — da arquitetura à aprovação em demo ao vivo com o cliente — habilitando terminais nos fluxos de digital bunkering.
            \item \textbf{Backend \& Integração:} Manutenção corretiva e desenvolvimento de features em API Java, garantindo integridade dos dados em operações marítimas internacionais.
            \item \textbf{Ambiente Internacional:} Colaboração diária em inglês com equipe em Singapura; suporte técnico direto a clientes internacionais e gestão de tickets via Mantis em rituais ágeis.
		\end{highlights}
	\end{onecolentry}

	\vspace{0.2 cm}

	\begin{twocolentry}
		{ Fev 2026 - Atual } \textbf{Analista de Automação}, Tyconnex / Schneider Electric (Singapura)
	\end{twocolentry}

	\vspace{0.10 cm}
	\begin{onecolentry}
		\begin{highlights}
			\item \textbf{Comissionamento FPSO (Petrobras):} Engenheiro de comissionamento \textit{in loco} dos sistemas AVEVA Plant SCADA (Citect) em duas plataformas offshore (P83 e P80) — 50.000+ tags, 1.700+ dispositivos de campo, PLCs Schneider M580.
			\item \textbf{IEC 61850 \& Protocolos:} Configuração e validação de relés de proteção SEL (F650, 869, 845) e VFDs Danfoss via IEC 61850 MMS e Modbus TCP; mapeamento de sinais RTU e integração com lógica ESA.
			\item \textbf{Integração EPO (Automotivo):} Comissionamento de relés de proteção SEL 2414 no AVEVA Power Operation (EPO) para planta automotiva no Brasil — configuração de dispositivos IEC 61850, endereçamento de tags e validação de polling.
			\item \textbf{Ferramenta de Campo:} Desenvolvimento de aplicação web offline Tag Finder (HTML único, 4MB) consolidando dados do projeto SCADA para diagnóstico em campo em tempo real — cobrindo 5.000+ equipamentos nas duas plataformas.
		\end{highlights}
	\end{onecolentry}

	\vspace{0.2 cm}

	\begin{twocolentry}
		{ Ago 2024 - Mai 2025 } \textbf{Desenvolvedor de Sistemas}, MAIMA Tecnologia
	\end{twocolentry}

	\vspace{0.10 cm}
	\begin{onecolentry}
		\begin{highlights}
			\item \textbf{Automação SAP (RPA):} Desenvolvimento de robôs em Python para automação de 4--6 processos críticos diários no SAP, reduzindo o tempo de execução manual de 8 horas para 2 minutos.
			\item \textbf{Manutenção Preditiva:} Sistema em C\# para previsão de paradas de máquina com base em análise de qualidade de produto (cerdas dentais), antecipando falhas críticas antes da ocorrência.
            \item \textbf{Arquitetura de Microsserviços:} Segregação em microsserviços distribuídos (Análise de Imagem, Dados, BFF), garantindo escalabilidade e manutenção simplificada.
		\end{highlights}
	\end{onecolentry}

	\vspace{0.2 cm}

	\begin{twocolentry}
		{ Jun 2022 - Ago 2024 } \textbf{Técnico em Automação}, S.Build Automação
	\end{twocolentry}

	\vspace{0.10 cm}
	\begin{onecolentry}
		\begin{highlights}
			\item \textbf{Engenharia de Campo \& Comissionamento:} Especialista em implementação de ecossistemas \textbf{AVEVA System Platform}, desenvolvimento de interfaces \textbf{InTouch} e configuração de \textbf{Historian} em ambientes industriais críticos.
            \item \textbf{Stack Industrial:} Projetos com PLCs Siemens (S7-1200/1500), Rockwell (Control/Compact/MicroLogix) e Schneider. Redes: Profinet, Ethernet/IP, Modbus, OPC. Workflows: \textbf{AVEVA WorkTasks}.
            \item \textbf{Automação de Dados \& Retrofit:} Scripts Python para atualização em massa de registros no AVEVA MES; migração de legado (1995 $\rightarrow$ 2023) e retrofit por engenharia reversa; upgrade de 20+ servidores via ThinManager.
		\end{highlights}
	\end{onecolentry}

    \vspace{0.2 cm}

    \begin{twocolentry}
		{ Set 2021 - Fev 2023 } \textbf{Competidor WorldSkills (Mecatrônica)}, SENAI Roberto Simonsen
	\end{twocolentry}

	\vspace{0.10 cm}
	\begin{onecolentry}
		\begin{highlights}
			\item \textbf{Campeão Estadual (São Paulo Skills):} Selecionado como \textbf{Melhor Aluno (Top 1)} do SENAI Roberto Simonsen para a equipe de elite; equipe sagrou-se \textbf{Campeã Estadual}. Rotina de 8--10h diárias de montagem de precisão, programação de PLCs e estações Festo MPS sob pressão de competição.
		\end{highlights}
	\end{onecolentry}

    \section{Certificações}

    \begin{highlights}
        \item \textbf{CSI AVEVA} Manufacturing Execution System (MES) 2020 - Operations \& Performance Exams
        \item \textbf{CSI AVEVA} Application Server 2023 Exam
        \item \textbf{Formação .NET Developer} (DIO - 93h) \quad \textbf{Java OOP} (Udemy - 54.5h) \quad \textbf{Node.js Fundamentals} (DIO - 30h)
    \end{highlights}

	\section{Educação}

	\begin{twocolentry}
		{2023 - 2025} \textbf{UNINTER}, Análise e Desenvolvimento de Sistemas
	\end{twocolentry}

	\begin{twocolentry}
		{2021 -  2022} \textbf{SENAI Roberto Simonsen}, Técnico em Redes de Computadores
	\end{twocolentry}

    \begin{twocolentry}
		{2020 -  2022} \textbf{SENAI Roberto Simonsen}, Eletricista de Manutenção Eletroeletrônica
	\end{twocolentry}

	\section{Idiomas}

	\begin{onecolentry}
		\begin{highlights}
			\item Inglês (C1 - Avançado)
			\item Português (Nativo)
		\end{highlights}
	\end{onecolentry}

\end{document}
